% Generated draft: requires mathematical and source review.
\documentclass[11pt]{article}
\usepackage[margin=25mm]{geometry}
\usepackage{fontspec}
\setmainfont{Latin Modern Roman}
\usepackage{amsmath,amssymb,mathtools,unicode-math}
\setmathfont{Latin Modern Math}
\usepackage{enumitem,needspace,xurl,hyperref}
\setlength{\parindent}{0pt}
\setlength{\parskip}{5pt}
\setlength{\emergencystretch}{3em}
\newcommand{\N}{\mathbb{N}}
\newcommand{\Z}{\mathbb{Z}}
\newcommand{\Q}{\mathbb{Q}}
\newcommand{\R}{\mathbb{R}}
\newcommand{\C}{\mathbb{C}}
\newcommand{\F}{\mathbb{F}}
\newcommand{\A}{\mathbb{A}}
\newcommand{\PP}{\mathbb{P}}
\newcommand{\E}{\mathbb{E}}
\newcommand{\eps}{\varepsilon}
\newcommand{\dd}{\,\mathrm{d}}
\begin{document}
\section*{UnsolvedMath ID: 800151}
\textbf{Source problem code:} catalogue entry
\subsection{Definitions and mathematical statement}
For a positive integer \(n\), let \(\Omega(n)\) be the number of prime factors of \(n\), counted with multiplicity, with \(\Omega(1)=0\). Define the real Dirichlet character modulo \(4\) by
\[
\chi_{-4}(n)=\begin{cases}0,&2\mid n,\\1,&n\equiv1\pmod4,\\-1,&n\equiv3\pmod4.\end{cases}
\]
For \(\varepsilon\in\{1,-1\}\), put
\[
S_\varepsilon=\{n\in\N:n\ge1,\ (-1)^{\Omega(n)}\chi_{-4}(n)=\varepsilon\}.
\]
Thus even integers belong to neither set. For a set \(S\) of positive integers, its logarithmic density, when it exists, is
\[
\delta_{\log}(S)=\lim_{x\to\infty}\frac{1}{\log x}\sum_{\substack{1\le n\le x\\n\in S}}\frac1n.
\]
\textbf{Requested calculation, attributed to Greg Martin.} Assuming GRH and LI, establish the existence and calculate the values of
\[
\delta_{\log}(S_1)\qquad\text{and}\qquad\delta_{\log}(S_{-1}).
\]
Here GRH refers to the generalized Riemann hypothesis for the relevant Dirichlet \(L\)-functions. LI refers to a linear independence hypothesis on positive imaginary parts of nontrivial zeros over \(\Q\). The supplied record does not specify the family of \(L\)-functions, the treatment of repeated ordinates, or the exact version of LI; these conditional hypotheses therefore require clarification for a fully specified formulation.

The historical motivation is attributed to Shanks (1959): for every integer \(k\ge1\), odd integers \(n\) satisfying \(\Omega(n)=k\) should be biased toward
\[
n\equiv(-1)^k\pmod4.
\]
The record identifies \(k=1\) with Chebyshev's original observation and describes the motivation as a bias of \((-1)^{\Omega(n)}\chi_{-4}(n)\) toward \(1\). It gives no precise definition of this fixed-\(k\) bias. The displayed density question concerns the pointwise values at integers \(n\); it does not specify a density of real parameters at which one counting function exceeds another.

The source also suggests, conjecturally, a similar phenomenon for \((-1)^{\Omega(n)}\chi(n)\) for any quadratic Dirichlet character \(\chi\), particularly the nonprincipal character modulo \(3\),
\[
\chi_{-3}(n)=\begin{cases}0,&3\mid n,\\1,&n\equiv1\pmod3,\\-1,&n\equiv2\pmod3.\end{cases}
\]
No precise additional density calculation, definition of the suggested bias, or restriction concerning principal or imprimitive characters is supplied.

\textbf{Additional passage attributed to Daniel R. Johnston.} For real \(t\ge0\), let \(\pi(t)\) count primes at most \(t\). For a positive integer \(q\) and an integer \(a\), let \(\pi(t;q,a)\) count primes \(p\le t\) with \(p\equiv a\pmod q\), and let \(\phi(q)\) be Euler's totient function. The supplied passage reports, citing \texttt{johnston2023average}, Theorem 1.1, the equivalence
\[
\mathrm{RH}\quad\Longleftrightarrow\quad
\mathfrak A_1^\pi(x):=\int_0^x\bigl(\pi(t)-\operatorname{li}(t)\bigr)\,\dd t<0
\quad\text{for every real }x>2.
\]
RH here denotes the Riemann hypothesis for the Riemann zeta function. The symbol \(\operatorname{li}\) denotes the logarithmic integral, but its normalization and the interpretation of the integral across its singularity are not specified in the record; these conventions must be fixed before treating the reported equivalence as an exact statement.

The passage introduces
\[
\mathfrak A_1^\pi(x;q,a):=\int_0^x\bigl(\phi(q)\pi(t;q,a)-\operatorname{li}(t)\bigr)\,\dd t.
\]
It explains the reported RH inequality heuristically through the size of a zero contribution involving \(\sum_\rho 1/(\rho(\rho+1))\), and attributes to Martin, in private communication, the observation that the analogous contribution from Dirichlet \(L\)-function zeros grows faster than the bias as \(q\) increases. It consequently suggests that \(\mathfrak A_1^\pi(x;q,a)<0\) should fail for large \(q\). This is supplied motivation, not a quantified additional problem: the record does not specify the admissible residue classes, the quantifiers on \(q,a,x\), the applicable hypotheses, or the precise zero sum and comparison. No such specifications are inferred here.
\subsection{Short English statement}
Assuming the source's GRH and LI hypotheses, determine the logarithmic densities of positive integers for which \((-1)^{\Omega(n)}\chi_{-4}(n)\) equals \(1\) or \(-1\). The precise LI hypothesis needs clarification. The record also suggests an extension to quadratic characters and appends an incompletely quantified discussion of integrated prime-counting errors in arithmetic progressions.
\subsection{Sources}
\begin{itemize}
\item Supplied local UnsolvedMath registry record: permanent ID \(800151\), source code \texttt{catalogue entry}, title ``Comparative Prime Number Theory --- Problem 20,'' in the AMR Open Problem Lists collection. The record labels the problem \emph{open}; this is a supplied status, not an independently verified mathematical or literature assessment.
\item Hamieh, Kadiri, Martin, and Ng, \emph{Comparative Prime Number Theory Problem List} (2024), Problem \(\#20\), \url{https://arxiv.org/abs/2407.03530}. The registry records access on 2026-10-16 and extraction by \texttt{arxiv\_latex}. These are supplied provenance details; the source has not been independently visited or verified for this draft.
\item The density question is attributed in the extracted text to Greg Martin, University of British Columbia. The record's list-level proposer field names Hamieh, Kadiri, Martin, and Ng, with proposed year 2024.
\item Shanks (1959), cited in the supplied text by the reference label \texttt{[Shanks]}, is credited with the fixed-\(k\) combinatorial observation. No full bibliographic entry for this reference is supplied.
\item The appended passage is attributed to Daniel R. Johnston, University of New South Wales--Canberra at ADFA. Its claimed RH equivalence cites \texttt{[johnston2023average, Theorem 1.1]}; no full bibliographic entry is supplied. The observation about arithmetic progressions is attributed to Martin by private communication.
\end{itemize}
\end{document}
